% 卷末回顾：由 main.tex 的无编号章引入，不改变正文层级与习题编号。
\begingroup
% 避免段落或列表项只留一行在相邻页。
\clubpenalty=10000
\widowpenalty=10000
本卷从安装 Ubuntu 开始，经过命令行、版本管理和现代 C++，最后走到了多文件工程的构建与运行。
这些内容共同回答一个问题：\textbf{拿到一份 C++ 工程，怎样在自己的电脑上读懂它、修改它、构建它，并把可验证的修改交给队友？}
卷首介绍的视觉组职责，也因此有了具体的落脚点：图像处理、目标识别与后续决策需要算法，而算法要成为可以协作开发、稳定运行的程序，还需要本卷这些基本功。

下面按三章回顾知识，再把它们放回同一条开发流程中。复习时可以一边读，一边打开终端或已有练习工程验证。
正文中已经给出的安装步骤、代码示例和参数表，仍是动手时的查阅入口；这里着重说明每部分解决什么问题、与其他内容怎样衔接，以及应该掌握到什么程度。

\subsubsection*{一、Linux 工作环境：让开发有一个可靠的起点}

\textbf{1. 安装系统：先认清磁盘，再建立工作环境。}
双系统安装依次涉及分区、制作启动盘、从 U 盘安装和安装后的系统设置。
分区决定两套系统各自使用哪些存储空间，启动盘提供安装环境，安装时的挂载点决定 Linux 如何使用这些分区。
复习时应能说明根目录 \texttt{/}、家目录 \texttt{/home}、交换空间与 EFI 系统分区各自的用途，并能对照自己的磁盘布局指出对应位置。
安装界面上的设备名、容量、文件系统和挂载点，要结合实际机器核对，不能照搬截图中的盘符与数字。

安装完成的判断也要落到操作上：Ubuntu 和 Windows 都能正常启动，网络与常用外设可用，能够打开终端、安装软件并保存自己的文件。
系统扩容、GRUB 引导菜单美化和桌面美化属于选修；前者服务于空间调整，后两者改善使用体验。
遇到这类需要时，应回到对应小节查具体步骤，并先确认当前系统与磁盘的状态。

\textbf{2. 备份与恢复：给环境修改留下退路。}
本章介绍了 Timeshift 系统快照和用 \texttt{tar} 打包的系统备份。
快照主要用于在系统配置出错后回到先前状态；归档备份是否完整，则取决于实际纳入的目录、排除规则和保存位置。
能够创建备份只是第一步，还应知道备份保存在哪里、包含哪些内容，以及恢复时需要什么环境。
个人代码、数据和重要文档也要有明确的备份安排，不能默认系统快照已经包含它们。

卸载重装一节进一步涉及系统分区与引导关系：移除 Ubuntu 的分区之后，还要考虑剩余系统能否正常引导以及磁盘空间如何重新使用。
这些维护内容按选修复习，但配置环境前先备份的习惯应保留。
能恢复的环境，才适合放心地继续实验。

\begin{note}
格式化、删除分区、系统恢复和覆盖文件都可能造成数据丢失。执行前应确认目标设备与目录，确认重要数据已有可用备份。
终端中的 \texttt{rm} 删除通常不会进入桌面回收站；不要通过反复执行破坏性命令来试探路径是否正确。
\end{note}

\textbf{3. 文件系统：用目录、路径和挂载组织资源。}
Linux 以 \texttt{/} 为起点组织目录树，绝对路径从根目录写起，相对路径则依赖当前工作目录。
\texttt{/home} 通常保存用户文件，\texttt{/etc} 保存许多系统配置，\texttt{/dev} 提供设备文件入口，\texttt{/proc} 暴露进程与内核的相关信息。
“一切皆文件”帮助我们用相近的接口访问不同资源，但普通文件、目录和设备文件仍有不同用途，不能随意互换操作。
挂载则是把某个文件系统接入目录树：找到挂载点，才能沿路径访问其中的文件。

路径概念会贯穿后两章。源码在哪里、程序从哪里读取图片、构建目录建在哪里，都与当前目录和路径写法有关。
遇到“文件不存在”，先核对路径、大小写和工作目录，往往比立刻修改代码更有效。

\textbf{4. 常用命令：从管理文件到观察系统。}
本章的命令可以按实际任务串起来：
\begin{itemize}
    \item \textbf{用户与权限}：区分普通用户和管理员操作，读懂文件的所有者、所属组与其他用户权限，使用 \texttt{chmod} 调整读、写、执行权限。权限不足时，先判断需要哪一项权限。
    \item \textbf{系统信息与软件安装}：用 \texttt{uname} 查看系统信息，用 \texttt{free} 观察内存，用 \texttt{df} 查看文件系统容量，用 \texttt{du} 统计目录占用。通过 \texttt{apt} 管理软件包时，要分清更新软件包索引、安装软件和升级已安装软件。
    \item \textbf{文件与目录管理}：用 \texttt{pwd}、\texttt{cd}、\texttt{ls} 定位，用 \texttt{mkdir}、\texttt{cp}、\texttt{mv}、\texttt{rm} 完成创建、复制、移动和删除，用 \texttt{tar}、\texttt{zip} 归档或压缩文件。操作前要能解释源路径与目标路径。
    \item \textbf{链接}：硬链接为同一文件提供另一个名字，软链接保存通向目标的路径。复习时应能解释为何目标被移走后软链接可能失效，以及为什么删除一个硬链接不等于立即删除所有文件内容。
    \item \textbf{查找与处理输出}：\texttt{find} 按文件名、类型等条件查找文件，\texttt{grep} 按模式筛选文本；管道把前一条命令的标准输出交给后一条命令，重定向把输入或输出接到文件等位置。要分清覆盖写入与追加写入。
    \item \textbf{阅读与编辑}：用 \texttt{cat}、\texttt{head}、\texttt{tail} 等查看内容，利用 Vim 的模式切换完成输入、保存与退出。阅读配置文件和构建日志，是这些操作最常见的用途。
\end{itemize}
AstroNvim 是本章介绍的选修编辑环境。它提供更完整的终端编辑体验，使用它时仍会用到文件定位、Vim 操作和配置管理这些基础。
常用命令应练到能够独立完成任务；不常用的参数，通过参数表和手册查阅即可。

\textbf{5. Shell 与环境变量：把一次操作变成可以重复的流程。}
终端提供交互界面，Shell 解释并执行命令；Bash 是本卷主要使用的 Shell。
脚本把一系列命令保存为文件，第一行的解释器声明用于直接执行脚本时选择解释器。
直接运行 \texttt{./test.sh} 需要执行权限；写成 \texttt{bash test.sh} 则是显式让 Bash 读取脚本。
应能把已经手动验证的准备目录、启动程序等操作整理成简单脚本，并解释每一步的效果。

环境变量影响进程如何寻找程序与读取配置。\texttt{PATH} 保存命令搜索目录，追加路径时要保留原有值。
在终端中设置变量、写入用户的 \texttt{.bashrc}、修改系统级配置，生效对象与时机不同。
还要分清“启动一个子进程执行脚本”和“用 \texttt{source} 在当前 Shell 中读取脚本”：前者对环境的修改不会自动带回父 Shell。
Zsh、Oh My Zsh 与提示符美化属于选修，切换 Shell 后也应检查使用的是哪份配置文件。

\textbf{6. Git：把修改变成有来由、可协作的版本。}
Git 的主线是工作区、暂存区、本地版本库与远程仓库之间的流转。
编辑文件改变工作区，\texttt{git add} 选择本次准备提交的内容，\texttt{git commit} 在本地形成版本，\texttt{git push} 把提交发送到远程。
\texttt{git clone} 用于取得仓库，\texttt{git pull} 用于获取并整合远程更新。
提交前查看状态与差异，能帮助我们确认提交的是预期修改；提交说明则帮助队友和未来的自己读懂修改目的。

提交者信息用于标记版本作者，HTTPS 令牌或 SSH 密钥用于远程访问认证，二者承担不同职责。
\texttt{.gitignore} 规定哪些未跟踪文件应被忽略，适合排除构建产物等本地文件，但不会自动取消对已跟踪文件的管理。
分支用于组织不同开发线路，合并与变基用于整合修改，子模块用于在工程中引用独立仓库。
日常提交与同步需要熟练；阅读分支历史、处理冲突和使用子模块时，要先认清当前状态，再查相应操作，避免凭命令名称猜测效果。

\textbf{7. Conda（选修）：把不同项目的依赖分开。}
Conda 一节介绍安装配置，以及环境的创建、激活、退出和删除。
它让不同项目能够使用各自的 Python 与软件包组合，减少依赖互相影响。
使用时要先确认当前激活的是哪个环境，再安装或运行软件；环境激活还可能改变 \texttt{PATH}，因此它与前面的环境变量知识直接相关。
后续进入深度学习等专项时，这种按项目管理环境的习惯会继续发挥作用。

\subsubsection*{二、现代 C++ 基础：组织数据、行为、资源和并发任务}

\textbf{1. 类与对象：把相关数据和操作放在一起。}
类描述一类对象具有哪些数据和行为，对象是具体实例。
数据成员保存状态，成员函数提供操作，构造函数建立初始状态，析构函数处理对象结束生命周期时的清理。
\texttt{public}、\texttt{private} 与 \texttt{protected} 用来划定访问边界：调用者使用公开接口，类内部维护自己的状态。
复习时可以拿装甲板对象举例，说明哪些数据需要保存、哪些操作应由成员函数完成，以及哪些状态不应让外部任意修改。

静态成员与类本身相关，友元为指定函数或类提供访问权限，\texttt{this} 指向当前对象。
函数重载通过不同参数列表提供同名操作，不能只靠返回类型区分；运算符重载则让自定义类型支持相应运算。
构造函数的初始化列表、拷贝构造和赋值也要分清：创建一个新对象，与修改一个已经存在的对象，是不同的时机。

\textbf{2. 继承、多态与模板：让接口和代码能够复用。}
继承建立基类与派生类的关系；虚函数允许通过基类指针或引用调用实际对象对应的实现。
\texttt{override} 帮助编译器检查重写关系，纯虚函数用于约定需要派生类完成的接口，抽象类不能直接实例化。
这里应能区分“同名函数的重载”和“派生类对虚函数的重写”，并通过调用结果验证自己的判断。

函数模板和类模板把类型等信息作为参数，使同一份代码适用于不同实例。
阅读容器、智能指针和通用函数时，都需要这种视角。
本卷要求能写简单模板、识别实例化与类型推导；复杂模板技巧留到实际需要时继续学习。

\textbf{3. Lambda：把一小段行为放到使用它的位置。}
Lambda 表达式适合排序比较器、回调和线程任务等场景。
捕获列表说明如何取得外部状态，参数列表说明调用者需要传入什么，函数体描述具体工作；必要时可用 \texttt{std::function} 保存符合指定调用形式的可调用对象。
泛型 Lambda、初始化捕获和 \texttt{this} 捕获，扩展了这种写法的适用范围。

值捕获保存副本，引用捕获访问原对象。
选择捕获方式时，除了考虑是否要修改原变量，更要考虑对象能活多久：若回调在线程或队列中稍后执行，被引用的局部变量可能已经销毁。
能够写出 Lambda 之后，还应能说清它依赖哪些外部对象，以及执行时这些对象是否仍然有效。

\textbf{4. 引用、拷贝与移动：说明函数怎样使用对象。}
引用为已有对象提供别名。按值传参会形成参数对象，引用传参可以直接访问原对象；\texttt{const T\&} 常用于只读访问并避免不必要的拷贝。
左值与右值描述表达式的类别，左值引用和右值引用的绑定规则则影响重载选择与资源转移。
尤其要记住：\textbf{有名字的右值引用变量，在表达式中使用这个名字时仍是左值。}

\texttt{std::move} 进行类型转换，为选择移动操作创造条件，本身不搬运数据；实际是否转移资源，要看随后调用的构造函数或赋值操作。
仅仅把一个引用绑定到转换后的表达式，并不等于原对象的资源已经被取走。
\texttt{std::forward} 用于模板转发，按照推导出的类型保留实参的值类别；它与无条件表达移动意图的 \texttt{std::move} 各有用途。
复习正文中的输出示例时，应逐次判断发生的是构造、拷贝、赋值、移动，还是单纯建立别名。

\textbf{5. 智能指针：把资源释放责任写进类型。}
资源获取即初始化（RAII）把资源的持有与对象生命周期联系起来；智能指针用这种方式管理动态对象，减少遗漏释放与重复释放的风险。
三种指针的关键区别在于所有权：
\begin{itemize}
    \item \texttt{unique\_ptr} 表达独占所有权，不能拷贝，可以移动。一个明确的持有者负责对象生命周期时，优先考虑它。
    \item \texttt{shared\_ptr} 表达共享所有权，通过引用计数管理对象；最后一个共享所有者释放后，对象才会销毁。
    \item \texttt{weak\_ptr} 观察共享对象而不增加强引用计数，使用前通过 \texttt{lock} 尝试取得有效的共享指针，适合打破循环引用。
\end{itemize}
通常通过 \texttt{make\_unique}、\texttt{make\_shared} 创建相应对象。
两个对象互相持有 \texttt{shared\_ptr}，可能让引用计数始终无法归零；把合适的一方改成不拥有资源的观察关系，才能解决所有权环。
另外，共享所有权不代表所指向的数据可以被多个线程任意同时修改，数据同步仍需要单独设计。

\textbf{6. STL：根据访问方式选择容器。}
标准模板库（STL）把常用数据结构、迭代器与算法组合起来。
选择容器之前，先考虑数据如何进入、如何查询、如何删除，以及是否需要顺序：
\begin{itemize}
    \item \textbf{顺序容器}：\texttt{array} 适合固定大小数据，\texttt{vector} 提供连续存储和按下标访问，\texttt{deque} 支持两端操作，\texttt{list} 等链式容器体现节点连接的组织方式。\texttt{vector} 的大小与容量不同，扩容可能重新分配存储。
    \item \textbf{容器适配器}：\texttt{queue} 对应先进先出，\texttt{stack} 对应后进先出，\texttt{priority\_queue} 每次提供当前优先级最高的元素。它们限制访问方式，不提供普通容器那样的迭代器遍历接口。
    \item \textbf{树、堆与映射}：树的层次结构、堆的优先级关系帮助解释相应数据结构的行为。\texttt{map} 保持键有序，\texttt{unordered\_map} 基于哈希组织键值对；查找复杂度、遍历顺序和实际开销都需要结合任务考虑。
\end{itemize}
按键访问时也要留意操作是否修改容器：例如映射的下标操作可能插入不存在的键，只查询时可使用 \texttt{find} 或 \texttt{count}。
应熟练操作 \texttt{vector}，能够为排队、撤销、优先处理和按编号查询等需求选择合适的容器。

\textbf{7. 迭代器与算法：在修改容器时保持访问有效。}
迭代器把容器与算法连接起来。常见遍历范围是左闭右开的 \texttt{[begin(), end())}，\texttt{end()} 表示尾后位置，不能解引用。
插入、删除或重新分配存储后，原有迭代器、指针或引用可能失效，具体规则取决于容器与操作。
例如 \texttt{vector} 重新分配存储后，指向原存储的迭代器都会失效；删除元素时，删除位置及其后的迭代器会受到影响。

遍历删除时可以用 \texttt{erase} 的返回值继续前进；批量删除时，要分清 \texttt{remove\_if} 调整逻辑范围与 \texttt{erase} 真正改变容器大小的作用。
排序时将 Lambda 用作比较器，就把本章的行为表达与容器操作连在了一起。
这部分的掌握标准是能够解释一次修改后哪些访问仍有效，而不只是记住某段循环写法。

\textbf{8. 多线程：让多个任务有序地共享资源。}
线程共享所属进程的资源，并发表示多个任务交错推进，并行则要求它们确实同时执行。
创建 \texttt{std::thread} 之后，应安排好线程的结束与回收；\texttt{join} 等待线程完成，\texttt{detach} 解除线程对象与执行线程的关联，后者尤其需要谨慎处理被访问对象的生命周期。
多个线程访问共享数据，只要涉及冲突的读写，就需要考虑同步。

互斥锁保护共享状态，\texttt{lock\_guard} 等 RAII 锁把解锁与作用域绑定。
共享锁允许多个读者同时进入，写者则需要独占；原子变量适合相应的单次原子操作，但多个原子操作组合起来不自动成为一个不可分割的整体。
条件变量用于等待某个状态发生变化，等待时应检查谓词，并由互斥锁保护相关共享状态。

\texttt{promise} 提交异步结果，\texttt{future} 获取结果或异常，\texttt{async} 提供按启动策略执行任务并取得结果的方式。
线程池进一步组合工作线程、任务队列、锁和条件变量，让已有线程反复处理任务；关闭时必须安排停止接收、唤醒等待者和回收线程等步骤。
这些工具分别解决互斥访问、等待条件和传递结果的问题，不能把“加一把锁”当成所有并发问题的统一答案。

生产者—消费者说明队列非空、非满时怎样协调，读者—写者说明共享访问与公平性的取舍，哲学家进餐说明多个资源之间可能形成循环等待。
死锁的四个必要条件是互斥、持有并等待、不可抢占和循环等待。统一获取顺序或使用合适的多锁获取工具，是本卷介绍的常用处理思路。
复习时应能在代码中指出共享状态、保护它的锁、等待条件，以及线程退出的路径。

\textbf{9. 内存与生命周期：为前面的错误建立解释。}
结合本章涉及的内存基础，还应建立常见进程内存布局的图景：代码与只读数据区域、静态存储数据、动态分配区域，以及线程调用栈。
栈帧保存一次函数调用需要的信息，局部对象的生命周期结束后，指向它的指针或引用可能悬垂；过深递归或过大的栈上对象可能耗尽栈空间。
动态对象的释放则取决于相应的资源管理机制。具体布局由实现与系统决定，不能仅凭变量名字判断它位于哪一块内存。
把这些知识与引用捕获、智能指针、容器扩容和线程退出联系起来，才能解释为什么有些程序“能编译，运行却不稳定”。

\subsubsection*{三、编译过程与 CMake：把源码组织成可以交付的程序}

\textbf{1. 编译过程：区分每一步的输入与产物。}
本章先用分步命令展示源码如何成为可执行文件：预处理处理包含与宏等指令，编译生成汇编代码，汇编生成目标文件，链接组合目标文件与所需库。
常见中间文件依次为 \texttt{.i}（C++ 常用 \texttt{.ii}）、\texttt{.s} 和 \texttt{.o}；\texttt{g++ -c} 停在生成目标文件之后，不进行最终链接。
声明告诉编译器接口是什么，定义提供实际实现；头文件被找到，只能说明相关声明等内容可见，并不代表对应实现已经参与链接。

静态库把目标文件组织成归档，链接器按需要取用其中内容；动态库则还涉及程序运行时的装载与查找。
因此，编译通过、链接通过和运行成功是不同检查点。理解它们的区别，才能正确解释后面的构建错误。

\textbf{2. CMake 工作流：配置、构建与安装各有职责。}
CMake 根据 \texttt{CMakeLists.txt} 描述生成构建系统，再由底层构建工具调用编译器与链接器。
\texttt{cmake -S/-B} 分别指定源码目录与构建目录，配置成功只表示构建规则已经生成；\texttt{cmake --build} 执行构建，\texttt{cmake --install} 按安装规则放置产物。
源码目录与构建目录分开，便于保留干净源码，也便于使用不同构建配置。

\texttt{cmake\_minimum\_required} 声明最低版本要求，\texttt{project} 声明项目及使用语言。
项目名与可执行目标名可以不同；Debug、Release 等构建类型影响相应配置，修改配置或源码之后还需要重新构建，才能运行更新后的程序。
缓存变量会保存在构建目录中，排错时也要检查正在使用哪一个构建目录和配置。

\textbf{3. 目标与依赖：明确生成什么、需要什么。}
\texttt{add\_executable} 创建可执行目标，\texttt{add\_library} 创建库目标；加入目标的源文件决定哪些实现参与构建。
\texttt{target\_include\_directories} 提供头文件搜索目录，\texttt{target\_link\_libraries} 描述链接依赖，编译特性与编译定义则分别表达语言要求和宏等配置。
目标名代表构建系统中的实体，变量展开代表读取变量内容，二者不能混用。

现代 CMake 把配置尽量附着到具体目标，再通过依赖关系传播使用要求。
\texttt{PRIVATE} 供当前目标使用，\texttt{PUBLIC} 供当前目标与下游使用，\texttt{INTERFACE} 供下游使用。
这些作用域描述构建配置的使用范围，与 C++ 类的成员访问权限并不是同一套机制。
应能对着一个库说明：哪些头文件是公开接口，哪些只供内部实现使用，使用这个库的程序需要继承哪些配置。

\textbf{4. 配置与安装：让构建描述能够重复使用。}
变量、列表、条件分支和选项用于表达可变化的构建配置；CMake 变量要通过相应命令传给编译器，才会影响 C++ 中的条件编译。
\texttt{function}、\texttt{include} 与模块搜索路径帮助拆分并复用 CMake 逻辑，打印变量可以检查配置阶段实际拿到了什么值。
这些补充内容的目标是读懂已有构建文件，并在需要时组织自己的配置。

安装规则决定产物的目录布局，安装前缀决定相对安装路径从哪里开始计算。
例如相对目的地 \texttt{bin} 会放在所选安装前缀下；构建目录中的程序与安装目录中的程序，需要分清正在运行的是哪一个。
安装步骤不能代替事先完成构建。

\textbf{5. 外部依赖与实例：把自己的代码接到现有库上。}
OpenCV 示例把查找依赖、创建目标、配置包含目录、链接库和运行验证串成了一次完整练习。
\texttt{find\_package} 查找已有包的信息，并不会自动替你安装依赖；自行安装 OpenCV 时，\texttt{OpenCV\_DIR} 应指向包含包配置文件的目录。
Python 中安装的 OpenCV 包不能直接替代这个 C++ 示例所需的开发包。

图像查看器既要验证正常输入，也要检查未传参数、图片不存在等情况；窗口显示还依赖可用的图形环境。
章末 ROS2 片段是选读材料，用于识别模块目录、库目标和依赖关系；它依赖原工程与 ROS2 环境，不能作为独立示例复制运行。
ament 的具体约定与图像算法分别留到卷二相应章节继续学习。

\textbf{6. 排错：先定位失败阶段，再检查对应对象。}
\begin{itemize}
    \item \textbf{配置失败}：检查工具、包是否可用，以及包配置文件路径与缓存设置。
    \item \textbf{编译失败}：检查首个有效诊断涉及的语法、类型、头文件与包含目录。
    \item \textbf{链接失败}：检查函数是否有定义、签名是否一致、实现文件或库是否参与链接，以及是否存在重复定义。
    \item \textbf{运行失败}：检查输入文件、工作目录、动态库搜索与运行环境；崩溃或结果异常还需要回到对象生命周期、边界访问与并发逻辑。
\end{itemize}
需要了解工具实际执行了什么时，可查看详细构建输出。
增量构建只更新受修改影响的产物，因此“没有重新编译所有文件”通常是正常行为；关键是能说明这次修改影响哪些目标，并运行验证结果。

\subsubsection*{四、把三章连起来：一次修改怎样成为可验证的成果}

假设现在接到一个练习任务：读取一组模拟目标数据，筛选并排序，最后输出结果。
可以沿着下面的顺序完成，并逐步增加复杂度：
\begin{enumerate}
    \item \textbf{准备环境与取得代码}：确认系统、工具和工作目录可用，用 Git 取得练习工程，查看说明与依赖要求，确认开始修改前的仓库状态。
    \item \textbf{设计数据与接口}：用类表达目标数据，用容器管理集合，用函数或 Lambda 表达筛选、排序规则；在接口处说明是读取、修改、复制还是转移对象。
    \item \textbf{明确生命周期}：判断普通对象与引用是否足够，需要动态资源时再选择合适的智能指针；确认回调引用的对象不会提前销毁。
    \item \textbf{组织与构建工程}：把声明、实现与入口按职责放入文件，使用 CMake 创建目标、加入源文件并设置依赖，完成配置与构建。
    \item \textbf{运行与定位问题}：先验证单线程的正常输入和边界情况，保存结果。确有并发需要时，再加入任务分工与同步，并检查线程能否正常退出。
    \item \textbf{整理与交付}：把重复操作写成脚本，记录依赖与构建运行方式，检查 Git 差异，提交源码、构建文件与必要说明，并按协作约定同步。
\end{enumerate}
Linux 管理文件与执行环境，C++ 决定程序行为，CMake 描述构建关系，Git 保存修改的历史。
这四部分共同支撑可复现的开发：换一个干净的构建目录，甚至换一位队友，仍能按说明得到一致的预期结果。

\subsubsection*{五、卷末自查：用已有工程证明自己掌握了什么}

各章习题用来检查具体知识点，卷末则适合带着已有练习工程做一次综合演示。
复习时可以逐项记录“能独立完成”“查阅后能完成”或“仍需重做”，把下一轮练习放在真正薄弱的位置。
\begin{enumerate}
    \item \textbf{环境与文件}：在 Ubuntu 中独立定位工程、创建与复制文件、调整脚本权限，解释常见目录与当前工作路径；能够查看磁盘和内存使用情况。
    \item \textbf{命令组合}：查找指定文件，从日志中筛选信息，将结果保存到文件，并解释管道与重定向各自处理了什么。
    \item \textbf{脚本与变量}：运行自己的准备脚本，说明解释器、执行权限和环境变量的生效范围；能判断新终端与当前终端的配置是否一致。
    \item \textbf{版本管理}：展示一次修改从差异检查、暂存到提交的过程，说明哪些文件应纳入版本库，哪些属于本地构建产物。
    \item \textbf{C++ 表达}：在已有代码中指出类的接口、构造与析构、模板和 Lambda，解释一次参数传递是否发生拷贝，以及移动操作是否真的被调用。
    \item \textbf{资源与容器}：解释一处对象由谁持有、何时销毁，说明容器的选择依据，并检查一段删除或插入代码是否继续使用了失效的迭代器。
    \item \textbf{并发协作}：在多线程练习中指出共享数据、同步方式、等待条件和退出流程，能够解释遗漏同步或改变加锁顺序会带来什么问题。
    \item \textbf{构建与复现}：在新的构建目录中重新配置、构建并运行多文件工程，按已有规则安装产物；展示一次错误的定位过程和修复后的验证结果。
\end{enumerate}
备份恢复、扩容、重装、桌面与终端美化、AstroNvim、Conda 等选修内容不计入通识考核，遇到实际需求时再回查与练习。
函数参数不必全部背下，复杂模板与进阶并发也需要后续积累；应当练熟的是常用操作，并能说清代码、资源与构建配置之间的关系。
参考题解放在本卷末尾，建议先独立作答或运行，再对照检查推理与操作遗漏。

\subsubsection*{六、带着这些基础进入专项学习}

完成本卷后，可以进入\textbf{卷二 · 专项技术栈}，按任务需要学习 ROS2、OpenCV、深度学习、估计与控制、导航等内容。
阅读 ROS2 节点时，会再次见到类、回调、智能指针和线程；处理图像时，会继续面对数据存储、对象共享与计算开销；接入外部库时，会继续使用依赖查找和目标链接。
自瞄、雷达与导航的具体任务不同，但都需要能运行、能排错、能协作修改的工程基础。

\textbf{卷三 · RM 视觉概述}再把技术放回比赛场景，讨论系统链路与工程准则。
学习专项时，可以随时带着实际问题回查本卷：找不到程序就查环境与路径，数据失效就查生命周期，构建失败就查目标和依赖。
当你能独立复现一个小工程、解释关键选择，并留下队友可以继续工作的代码与说明时，本卷的学习目标就落到了实处。
\par
\endgroup
